% 来源：OEBPS/Text/01_56.xhtml（OPF spine 顺序；仅第一册）
\cleardoublepage
\BookEntry[孔子《论语》]{孔子《论语》}

\BookCover{images/书单39论语.jpg}{孔子著| 中华书局2016年版}

\begin{BookSummaryBox}
\textbf{中国历史上最伟大的思想家孔子的语录体散文集}
\par
\end{BookSummaryBox}

\subsubsection*{这本书谈什么？}

《论语》是一部语录体思想散文集，来源于孔子给弟子讲学的内容。早在春秋后期孔子设坛讲学时期，其主体内容就已初始创成；孔子去世以后，他的弟子和再传弟子代代传授他的言论，并逐渐将这些口头记诵的语录言行记录下来，最终在战国初年汇辑论纂成书。
\par

孔子思想以立身为出发点，而人能立身于世的首要条件就是具有“君子”人格。君子具备仁爱之心，自重自律；表里如一，言行一致；积极进取，德才兼备；孜孜于学，注重实践；安贫乐道，谨守正义等等；始终坚守人之间的相处之道，遵从不同的伦理关系，构建着和谐友爱的人际关系和社会环境。
\par

\subsubsection*{芒格为什么谈起这本书？}

2007年5月23日，芒格在南加州大学的演讲中提到道：
\par

\begin{quotation}
我这辈子一直很崇拜孔子。我很喜欢孔子关于“孝道”的思想，他认为孝道既是天生的，也需要教育，应该代代相传。你们大家可别小看这些思想，请留意在美国社会中亚洲人的地位上升得有多快。我认为这些思想很重要。
\par

在孔子的思想中，孝道——对父母及祖先的敬爱——是一种需要培养的美德。广义的“孝道”包括照顾父母，听父母的话，敬爱和赡养父母，对父母有礼貌，传宗接代，兄弟之间相互扶持，给父母明智的建议，隐讳他们的错误，在他们生病和死亡时表达哀伤，在他们去世后进行祭祀等。
\par

孔子认为，如果人们能够学会践行孝道，他们就能够更好地扮演他们在社会和政府中的角色。在孔子看来，孝道十分重要，甚至比法律还重要。实际上，在汉朝的某些时期，那些不按孝道祭祀祖先的人会遭到肉体上的刑罚。
\par
\end{quotation}

\subsubsection*{子曰}

“成事不说，遂事不谏，既往不咎。”
\par

“智者上善若水，海纳百川；仁者高山仰止，厚德载物。”
\par

“君子坦荡荡，小人长戚戚。”
\par

“益者三友。友直，友谅，友多闻，益矣。”
\par

“吾十有五而志于学，三十而立，四十而不惑，五十而知天命，六十而耳顺，七十而从心所欲不逾矩。”
\par

“质胜文则野，文胜质则史，文质彬彬，然后君子。”
\par

“不愤不启，不悱不发，举一隅，不以三隅反，则不复也。”
\par

“君子和而不同，小人同而不和。”
\par

“饭疏食，饮水，曲肱而枕之，乐亦在其中矣。不义而富且贵，于我如浮云。”
\par

“弟子入则孝，出则悌，谨而信，泛爱众而亲仁，行有余力，则以学文。”
\par

“德之不修，学之不讲，闻义不能徙，不善不能改，是吾忧也。”
\par
